\mode*
\section{Search and verification protocol}
\label{app:search}

Every citation in this paper follows a find--verify--record protocol: the
claim is stated first, a source is found through a recorded, reproducible
query, the source itself is read (abstract or full text) to verify that it
supports the claim at its scope and strength, and the outcome is recorded
in a provenance block in the bibliography file (fields: claim, found-via,
picked, verbatim quote, verification status, date).
This appendix documents the searches; the per-reference records are in the
provenance blocks.
% XXX Future searches: prefer the scholar tool, which queries several
% databases at once and records the queries; the API queries below are
% recorded verbatim and equally reproducible.

\paragraph{Foundational references (2026-07-02)}

The approaches-to-learning and debugging-education anchors were found via
the Crossref REST API (\texttt{query.bibliographic}), one query each:
\begin{itemize}
  \item \texttt{on qualitative differences in learning outcome and process
    Marton} --- retained Marton and Säljö I and II (both parts surfaced in
    one query);
  \item \texttt{revised two-factor study process questionnaire R-SPQ-2F
    Biggs} --- retained Biggs, Kember and Leung;
  \item \texttt{Debugging a review of the literature from an educational
    perspective McCauley} --- retained McCauley et al.
\end{itemize}
Verification: abstracts from the Crossref records; the R-SPQ-2F was later
verified in full text (2026-07-09), from which the 20 items in
\cref{approach-measure} are transcribed.
(Semantic Scholar and OpenAlex were unavailable on 2026-07-02; Crossref
was used as the fallback.)

\paragraph{Passages in \citetitle{NCOL} (2026-07-02 to 2026-07-09)}

The claims attributed to \citeauthor{NCOL} were located by full-text
search of the book's chapter PDFs (\texttt{pdftotext} followed by
\texttt{grep}), with search terms \texttt{deep approach},
\texttt{induction}, \texttt{outcomes of learning turned out}, and
\texttt{point out} (the assessment discussion in chapter~4); page
numbers were taken from the printed running heads.
The studies Marton reviews (Silén; Fyrenius, Wirell and Silén; Marton, Wen
and Wong; Marton and Wenestam; Holmqvist, Gustavsson and Wernberg) were
identified from those passages; their bibliographic metadata was verified
against the book's reference list and, where indexed, against Crossref.
They are cited as reviewed in \citetitle{NCOL}; the originals remain to be
read.

\paragraph{Concept inventory for the diagnostic quiz (2026-07-09)}

Crossref \texttt{query.bibliographic} queries:
\texttt{language independent assessment CS1 knowledge FCS1 Tew Guzdial};
\texttt{replication validation language independent CS1 knowledge
assessment SCS1 Parker};
\texttt{revalidating cognitive introductory computing instruments
Bockmon}.
The IRT evaluation was missed by Crossref and found by the OpenAlex search
\texttt{item response theory evaluation SCS1}, which also surfaced the
FCS1/SCS1 retrospective.
Retained: the FCS1 and SCS1 primary sources, the IRT evaluation (later
read in full text for the item-level results), the shortened revalidation,
and the retrospective.
The Nordic Prior Knowledge Test was contributed by the first author
(2026-07-10); its metadata was verified against the publisher's page and
its full text read.

\paragraph{Related work (2026-07-10)}

Crossref \texttt{query.bibliographic} queries, verbatim:
\begin{itemize}
  \item \texttt{Debugging from the student perspective Fitzgerald
    McCauley} --- retained Fitzgerald et al.;
  \item \texttt{Improving debugging skills in the classroom systematic
    self-efficacy Michaeli Romeike} --- retained Michaeli and Romeike;
  \item \texttt{methods and tools for exploring novice compilation
    behaviour Jadud} --- retained Jadud;
  \item \texttt{ProgSnap2 flexible format programming process data Price}
    --- retained Price et al.;
  \item \texttt{Learning to program phenomenographic perspective Booth}
    --- Booth's thesis did not surface; an OpenAlex search with the same
    terms surfaced Bruce et al., which was retained instead (full text
    read from the journal's site);
  \item \texttt{Variation theory applied students learning computer
    programming Thune Eckerdal} --- retained Thuné and Eckerdal;
  \item \texttt{programming success factors deep surface approach Bergin
    Reilly} --- the candidate could not be verified (no retrievable
    abstract; truncated title metadata) and was \emph{discarded} rather
    than cited unread.
\end{itemize}
All retained related-work references were verified by reading their
abstracts (via OpenAlex or Semantic Scholar), except Bruce et al., whose
full text was read.

\paragraph{Criticisms of the R-SPQ-2F and the expert debugging literature
(2026-07-13)}

Searches via the \texttt{scholar} tool (which queries several databases at
once), verbatim:
\begin{itemize}
  \item \texttt{scholar search "Why Programs Fail a guide to systematic
    debugging Zeller" -p s2 -p dblp} --- retained the second edition of
    Zeller's book; verified against the table of contents and chapter
    descriptions on the book's companion site;
  \item \texttt{scholar search "R-SPQ-2F factor structure re-examination
    validity study process questionnaire" -p s2 -p openalex} --- surfaced
    application and translation studies but not the item-level
    re-examinations;
  \item \texttt{scholar search "Revised Two-Factor Study Process
    Questionnaire exploratory confirmatory factor analyses item level
    Justicia" -p openalex} --- surfaced Stes, De Maeyer and Van Petegem
    (retained; abstract and conclusions read on the journal's site) but
    still not Justicia et al., which was then found by the Crossref
    fallback \texttt{query.bibliographic} query \texttt{Revised Two-Factor
    Study Process Questionnaire exploratory confirmatory factor analyses at
    item level} (retained; abstract read on the publisher's page);
  \item \texttt{Exploring and reconciling the factor structure Revised
    Two-factor Study Process Questionnaire} (Crossref) --- surfaced Socha
    and Sigler's factor-structure study, but its abstract could not be
    retrieved from any source (publisher access only), so it was
    \emph{discarded} rather than cited unread.
\end{itemize}
The motive/strategy subscale composition and scoring key cited in the
criticisms discussion were verified against the full text of
\textcite{RSPQ2F} (p.~149).

\paragraph{Originals of the studies reviewed in NCOL (2026-07-19)}

The four studies backing the deep-approach reinterpretation
(\cref{debugging-as-learning}) had been cited as reviewed in
\citetitle{NCOL}; this round sought their full texts.
\begin{itemize}
  \item Web search \texttt{Silén "Mellan kaos och kosmos" avhandling
    2000 Linköping pdf} --- full thesis found in DiVA
    (diva2:24135) and read; the conclusion NCOL translates was
    located on p.~265, and the study's population corrected (the
    nursing programme, not medical students as the review states).
  \item Web search \texttt{Fyrenius Wirell Silén "approaches to
    achieving understanding" 2007 pdf diva} --- the journal article
    is paywalled; the author's doctoral kappa (Fyrenius,
    \emph{Dynamiskt lärande}, 2007, diva2:23720), found by listing
    her publications in DiVA, summarises the paper (its
    \emph{delarbete}~3) and was read instead.
  \item \Textcite{MartonWenWong2005}: Springer, S2 and OpenAlex all
    report closed access with no repository full text and no
    abstract on record; a web search for stray PDFs found none.
    \emph{Remains cited as reviewed in NCOL.}
  \item \Textcite{MartonWenestam1988}: the Wiley volume is in
    neither archive.org nor Google Books; no PDF surfaced by web
    search.  A page discrepancy between reference lists was noted
    for resolution against the printed book.
    \emph{Remains cited as reviewed in NCOL.}
\end{itemize}

\paragraph{Full texts read}

\Textcite{RSPQ2F} (instrument transcription),
\textcite{SCS1IRT} (item-level results),
\textcite{BruceLearningToProgram},
\textcite{NPKT},
\textcite{LiterateProgramming} (a known primary source, not searched
for; metadata verified against Crossref),
\textcite{Silen2000} (the thesis),
\textcite{FyreniusWirellSilen2007} (via the author's doctoral kappa,
see above),
and the passages of \citetitle{NCOL} cited throughout.
All other references are verified at abstract level, with
\textcite{DebuggingReview} cited only as a pointer to the review it is.

\paragraph{Grading and knowledge in Swedish independent schools (2026-08-24)}

This round backs the background claim that Swedish \emph{friskolor}
(independent, voucher-funded schools) award higher grades than
\emph{kommunala} (municipal) schools while their students know less.
The claim splits into two sub-claims that the literature treats very
differently, and the round was run to keep them apart:
(a)~friskolor grade more leniently \emph{relative to externally measured
knowledge}; and (b)~friskola students have \emph{lower knowledge} in
absolute terms.
The research question is therefore: \emph{is it true that Swedish
independent schools grade more generously than municipal schools, and is
it true that their students learn less?}

\subparagraph{Method}
Searches were run with the \texttt{scholar} tool.  A provider health check
(\texttt{scholar providers check}) was run first and is part of the record:
Semantic Scholar was rejecting its key (HTTP 403) and OpenAlex was rate
limited (HTTP 429) on the day, so the sweep rests on Web of Science and
Scopus, which are in any case the stronger indexes for
economics-of-education.  Verbatim queries:
\begin{itemize}
  \item \texttt{grade inflation independent schools voucher Sweden}
    (\texttt{-p wos -p scopus});
  \item \texttt{school competition grading standards national tests
    leniency} (\texttt{-p wos -p scopus});
  \item \texttt{TS=("grade inflation" AND (Sweden OR Swedish) AND (school
    OR voucher OR "independent school"))} (\texttt{-p wos});
  \item \texttt{TS=((voucher OR "independent school" OR "free school") AND
    (Sweden OR Swedish) AND (achievement OR "test score" OR grades))}
    (\texttt{-p wos}).
\end{itemize}
Counter-searches, phrased to surface work refuting or qualifying the
claim:
\begin{itemize}
  \item \texttt{voucher reform private schools long run effects student
    achievement Sweden} (\texttt{-p wos -p scopus});
  \item \texttt{independent school attendance effect achievement selection
    bias no difference} (\texttt{-p wos -p scopus});
  \item \texttt{TS=((Sweden OR Swedish) AND ("independent school" OR
    voucher OR "free school") AND (TIMSS OR PISA OR "international
    assessment" OR "externally graded"))} (\texttt{-p wos});
  \item \texttt{TS=((Sweden OR Swedish) AND school AND ("grading
    standards" OR "grade inflation") AND (criticism OR "no difference" OR
    overestimate OR contested))} (\texttt{-p wos}) --- \emph{zero hits};
  \item \texttt{TS=((Sweden OR Swedish) AND ("independent school" OR
    friskola OR voucher) AND ("higher achievement" OR outperform OR
    "positive effect") AND (school OR student))} (\texttt{-p wos}) ---
    \emph{zero hits}.
\end{itemize}
The two broad Scopus counter-queries returned 120 and 277 hits
respectively, almost all of them monographs on unrelated policy topics;
they were screened by title and venue and discarded, and the field-syntax
queries above were substituted as the reproducible record.
Full texts were read from the freely available working-paper versions
(IFAU, IFN, IZA) and checked against the published abstracts, since
ScienceDirect returns HTTP 403 to automated fetches.
Skolverket's reports are not indexed by any of the providers; they were
found by web search for
\texttt{Skolverket analyser likvärdig betygssättning huvudman fristående
kommunala skolor nationella prov betyg avvikelse rapport} and retrieved
from the agency, after \textcite{EdmarkPersson2021} pointed to them.

\subparagraph{Results}
Sub-claim~(a) is well supported, and by the one design that can establish
it.  \Textcite{HinnerichVlachos2017} observe the \emph{same} standardized
tests graded both internally by the school and externally by the Schools'
Inspectorate, and find voucher schools \(0.14\) standard deviations more
generous in their internal grading.  It is corroborated at both school
levels by Skolverket's full-population analyses, of year~9
\autocite{Skolverket2019Betyg} and of upper secondary
\autocite{Skolverket2020Betyg}; by \textcite{EdmarkPersson2021}, who find independent
students more often \enquote{up-graded} relative to their test result but
no more often down-graded, and historically by
\textcite{WikstromWikstrom2005}.
Sub-claim~(b) is \emph{not} supported.  In
\textcite{HinnerichVlachos2017} the raw difference in externally graded
test scores is \(-0.064\) standard deviations and not statistically
significant (Table~2, column~1 of the IFAU working paper); only the
value-added estimate---an effect of attendance conditional on prior
achievement, not a statement about absolute knowledge---reaches
significance, at \(-0.056\).  \Textcite{EdmarkPersson2021} conclude that
the added value of independent schools may be \enquote{close to zero},
not negative.  \Textcite{BohlmarkLindahl2015} run the opposite way
outright, finding that the expansion of the independent sector improved
average outcomes, primarily through external effects such as competition,
and finding no differential grade inflation at the municipal level.
Retained: all six references above.
Discarded after screening the abstracts: Ding, Klapp and Yang Hansen
(2024), whose comparison of public and independent schools turned out to
be about mathematics self-concept and self-efficacy rather than any
achievement gap; and IFAU working paper 2024:17 by Edmark, Hussain and
Haelermans, which concerns track choice and long-run labour outcomes.
Both bear on the Swedish voucher system but not on this claim.

\subparagraph{The mechanism was searched for and not found}
Because the intended use of the claim is to relate a student's school
history to their approach to learning
(\cref{approach-measure}), two searches looked specifically for
literature connecting school type to deep and surface approaches:
\texttt{TS=(("approach to learning" OR "surface approach" OR "deep
approach" OR "study process questionnaire") AND (school AND (private OR
independent OR voucher OR "school type" OR marketization)))}
(\texttt{-p wos}, 38 hits, all in medical, music or teacher education and
none about school type in this sense), and
\texttt{TS=(("approaches to learning" OR "deep learning approach" OR
"surface learning") AND (secondary OR "upper secondary" OR "high school")
AND (grades OR "grade inflation" OR assessment) AND (Sweden OR Swedish OR
Nordic))} (\texttt{-p wos}, one hit, on IQ and personality).
No study relates Swedish \emph{huvudman} type to the deep or surface
approach.  The link is ours to hypothesise, and is reported as such.

\subparagraph{Conclusion}
Yes on the first count, no on the second.  Independent schools do grade
more generously than municipal schools relative to externally measured
knowledge; the literature does \emph{not} show that their students know
less.  Three qualifications survive into the prose.
First, the leniency evidence is strongest at the upper-secondary level:
\textcite{BohlmarkLindahl2015} find no differential grade inflation at
the compulsory level, and \textcite{EdmarkPersson2021} suggest the
reconciliation that the problem may be specific to upper secondary.
Second, the effect is small and the \emph{huvudman} variable is a weak
predictor: Skolverket puts the share of between-school variation in
grading explained by huvudman at about \(1.5\)~per cent in year~9
\autocite{Skolverket2019Betyg} and \(0.5\)--\(1\)~per cent in upper
secondary \autocite{Skolverket2020Betyg}, with a \emph{wider} spread within
the independent sector than within the municipal one; relative grading
against a school's own performance level is a far larger factor.
Third, \textcite{HinnerichVlachos2017} explain their own findings by
schools responding to differences in educational \emph{preferences}, not
by any difference in how their students go about learning.
The safe formulation is therefore about grades and their relation to
demonstrated knowledge, never about friskola students knowing less.

\paragraph{Higher-education outcomes of friskola graduates (2026-08-24,
sub-round)}

A follow-up to the previous round asked whether \enquote{students from
friskolor dropped out to a higher extent than students from municipal
schools}.
The claim has two plausible readings, and they were checked separately:
(A)~drop-out from \emph{higher education} among students who came from an
independent \emph{gymnasium}, and (B)~drop-out \emph{within} gymnasiet by
school type.
The research question is: \emph{do students from Swedish independent
upper-secondary schools leave education at a higher rate than students
from municipal schools, and at which stage?}

\subparagraph{Method}
Reading~(A)'s source was reached by \emph{citation chaining} rather than
keyword search: \textcite{Skolverket2020Betyg} cites (p.~46, fn.~83) an
analysis showing that students from independent gymnasier perform worse
in their first year of higher education given the same upper-secondary
grades.  That reference was read off its reference list (p.~50) and the
report retrieved by the web search \texttt{Skolverket rapport 466 "Från
gymnasieskola till högskola" registerstudie 2018 pdf}, which also
surfaced the joint Skolverket/UKÄ successor report.  Both were read in
full.
Topic-driven queries were run to find any peer-reviewed treatment, and
returned almost nothing:
\begin{itemize}
  \item \texttt{TS=((Sweden OR Swedish) AND ("grade inflation" OR
    "inflated grades" OR "grading standards") AND (university OR "higher
    education" OR college) AND (performance OR dropout OR completion OR
    credits))} (\texttt{-p wos}) --- one hit, on gender discrimination in
    grading, not relevant;
  \item \texttt{TITLE-ABS-KEY((Sweden OR Swedish) AND ("upper secondary"
    OR gymnasium) AND ("independent school" OR voucher OR "free school")
    AND ("higher education" OR university) AND (performance OR dropout OR
    completion))} (\texttt{-p scopus}) --- one hit, a country overview,
    not relevant.
\end{itemize}
For reading~(B):
\texttt{TS=((Sweden OR Swedish) AND (dropout OR "drop out" OR attrition
OR completion OR "graduation rate") AND ("upper secondary" OR gymnasium
OR "high school") AND ("independent school" OR voucher OR private OR
"free school"))} (\texttt{-p wos}) returned a single hit, an essay on
equivalence and choice, with no completion analysis.  Skolverket's
statistics reports \emph{Betyg och studieresultat i gymnasieskolan}
(2025) and \emph{Enskilda huvudmän i grund- och gymnasieskolan}
(2024:4) were retrieved and searched; neither breaks completion down by
huvudman.
Counter-search: \texttt{TS=(("private school" OR "independent school" OR
voucher OR charter) AND ("college performance" OR "university
performance" OR "postsecondary performance" OR "college persistence" OR
"college GPA") AND (grades OR grading OR "grade inflation"))}
(\texttt{-p wos}) --- \emph{zero hits}.

\subparagraph{Results}
Reading~(A) is supported, but not as a claim about dropping out.
\Textcite{Skolverket2018Hogskola} finds that students from independent
gymnasier complete a smaller share of their first-year higher-education
credits than students from municipal schools, despite entering with
\emph{higher} upper-secondary grades, and that the gap survives controls
for sex, background, parental education, grades and choice of study
(odds ratios \(1.58\)--\(2.07\) for the risk of completing under half the
registered credits, all significant at the \(0.1\) per cent level).
\Textcite{UKASkolverket2024} replicates this across eight cohorts and
splits the independent sector: the shortfall is \(8.0\) percentage points
for for-profit and \(6.0\) for non-profit independent schools
(2017/18 cohort), and the report rules out the obvious artefact that
these students merely register for more credits.
Reading~(B) is unsupported: no study of gymnasium completion by school
type surfaced in any database, and the agency statistics do not report
it.
Retained: both reports.

\subparagraph{The wording matters}
Neither report measures drop-out.  The outcome is
\emph{prestationsgrad}, the share of \emph{registered} credits completed
in the first year, and both reports say explicitly that study
interruptions contribute to it without being separable from
it\autocite[p.~9]{Skolverket2018Hogskola}\autocite[p.~20]{UKASkolverket2024};
the 2018 study excludes students who left after one term from its
population altogether.
The literature therefore backs \emph{underperformance conditional on
grades}, not a higher drop-out rate, and the paper must say the former.

\subparagraph{Conclusion}
Yes at the transition to higher education, no within gymnasiet, and not
as \enquote{dropping out} in either case.  Two qualifications carry into
the prose.
First, the effect is \emph{conditional}: it is precisely because these
students arrive with higher grades and then complete less that the
finding is interesting, and it is evidence about what the grades measure
rather than about the students.  \Textcite{UKASkolverket2024} draws the
same inference (p.~47): grades \enquote{inte nödvändigtvis återspeglar
elevernas kunskapsnivåer}.
Second, at the \emph{extensive} margin the sign reverses:
\textcite{EdmarkPersson2021} find independent-school students more likely
to reach higher education and more likely to take any university credits
at all.  Independent-school students are thus more likely to get to
university and less likely to keep up once there; citing only one of
these would misrepresent the literature.
Both sources are agency register studies rather than peer-reviewed
research --- the topic-driven searches found no peer-reviewed treatment
in either direction --- and neither identifies a causal effect.
